% Cible : ARIEL_LOGMINER_MEMOIRE_FINAL/chapters/chapitre7_reproductibilite_deploiement.tex
% Passage à remplacer : fichier entier, de « \chapter{Reproductibilité » à la fin
% Agent : 6 — Reproductibilité, artefacts et déclaration IA
% Sources : chapitre 7 actuel ; annexes.tex, tab:preuves-experimentales ;
%           experiments/phase_revision_20261104/{README.md,MANIFEST_SHA256.csv} ;
%           HANDOFF.md §3 et §6 ; RETOURS_LECTEURS.md, lecteur 1 E.2.

\chapter{Reproductibilité, déploiement et exploitation}
\label{chap:reproductibilite}

\section{Portée de la reproductibilité}

La traçabilité expérimentale repose sur quatre éléments : une entrée identifiée, un
protocole paramétré, un identifiant de campagne et un artefact de sortie. La matrice de
l'annexe~\ref{tab:preuves-experimentales} relie les résultats principaux à leurs preuves.
Les fichiers CSV et JSON constituent la référence pour vérifier ou recalculer une valeur ;
les tableaux et les figures du manuscrit en sont des représentations.

Trois niveaux doivent être distingués. La \emph{traçabilité} permet de retrouver l'origine
d'une valeur. La \emph{répétabilité} consiste à réexécuter un script dans un environnement
équivalent. La \emph{reproductibilité indépendante} suppose qu'un tiers reconstruise le
résultat à partir du code, des données et des instructions remis. Le travail documente les
deux premiers niveaux ; aucune reproduction indépendante complète par un tiers n'est
revendiquée.

\section{Paquet de preuves et intégrité}

Les artefacts expérimentaux sont organisés par phase dans le répertoire
\texttt{experiments}. Le paquet
\texttt{experiments/phase\_revision\_20261104} regroupe les sorties de consolidation
utilisées pour HDFS, BGL, CSE-CIC-IDS2018, le routeur open-set, la corrélation, la
distribution multi-VM et le contrôle Isolation Forest. Son fichier
\texttt{MANIFEST\_SHA256.csv} recense 54 fichiers avec leur chemin relatif, leur taille et
leur empreinte SHA-256. Une vérification du manifeste établit l'intégrité binaire du
paquet ; elle ne prouve ni la provenance officielle des jeux de données ni la validité du
protocole.

Les données de travail volumineuses ou soumises à des conditions de redistribution ne
sont pas toutes incluses dans ce paquet. Leur provenance, leur version et, lorsqu'elle est
disponible, leur empreinte doivent donc accompagner toute reproduction. Les identifiants
d'accès, mots de passe et secrets de machines virtuelles sont exclus des preuves et doivent
être fournis par un gestionnaire de secrets ou des variables d'environnement.

\section{Procédure de reproduction}

Pour chaque campagne, la procédure minimale est la suivante :

\begin{enumerate}
  \item obtenir les données depuis la source documentée et vérifier leur empreinte lorsqu'une
  valeur de référence est disponible ;
  \item préparer l'environnement logiciel prévu pour la campagne et consigner les versions ;
  \item exécuter le script avec les paramètres, partitions et graines documentés ;
  \item contrôler le schéma des artefacts produits, puis comparer les métriques aux fichiers
  de référence ;
  \item vérifier les empreintes du paquet et seulement ensuite reconstruire les tableaux ou
  les figures.
\end{enumerate}

Une divergence doit être rapportée avec l'entrée, l'environnement, les paramètres et
l'artefact concernés. La régénération d'une figure ne constitue pas, à elle seule, une
reproduction du calcul qui l'alimente.

\section{Paramètres déterminants}

Les paramètres HDFS sont rassemblés dans le
tableau~\ref{tab:hdfs-parametres-annexe}. L'ordre des blocs selon leur première apparition,
la partition chronologique, l'échantillon d'ajustement de Drain3, l'état gelé du parseur,
l'agrégateur et le seuil choisi sur validation doivent être conservés. Une comparaison
HDFS n'est valide qu'au même niveau d'analyse : événement ou bloc.

Pour BGL, l'unité du protocole principal est l'événement. L'état Drain3 et les histogrammes
sont appris sur l'entraînement puis gelés ; le seuil est choisi sur validation. La position
temporelle des segments d'entraînement et de test est déterminante, car elle modifie la
couverture des templates et des sévérités.

Pour CICIDS2017 et CSE-CIC-IDS2018, la partition, la journée ou le scénario tenu à l'écart,
la prévalence du test et la règle de calibration du seuil doivent être publiés. Des scores
issus d'une séparation aléatoire, d'un holdout par scénario ou d'un transfert entre journées
ne répondent pas à la même question et ne peuvent pas être comparés sans cette information.

Le routeur open-set doit être rejoué avec la politique gelée, sans recalibrer son seuil sur
les sources évaluées. La corrélation doit publier la génération des incidents synthétiques,
les graines et la définition des paires positives et négatives. Ces conditions déterminent
directement les taux rapportés.

\section{Environnement multi-agent}

Le benchmark A--D doit conserver les architectures comparées, les niveaux de charge, le
nombre de répétitions et les unités de mesure. Les micro-latences dépendent de la machine ;
une reproduction sur un autre hôte doit publier sa configuration et vérifier les tendances
relatives plutôt que rechercher une égalité à la décimale.

La campagne multi-VM combine un coordinateur, une instance Redis centrale et deux agents
dont les profils sont décrits au tableau~\ref{tab:multivm-parametres-annexe}. Une
reproduction doit conserver la configuration, les versions des systèmes, les métriques par
tâche, les journaux des agents et le scénario de panne. La campagne HDFS distribuée évalue
le transport de scores déjà calculés et l'application d'un seuil ; elle ne reproduit ni le
parsing Drain3 ni le calcul des scores sur les machines virtuelles.

Redis demeure un point central du laboratoire. Une reproduction en réseau doit ajouter
authentification, contrôle d'accès et protection du transport ; une instance exposée sans
ces mécanismes ne correspond pas à un déploiement recommandé.

\section{Disponibilité des éléments et vérification externe}

Le dépôt remis avec le mémoire doit contenir le code du prototype, les scripts ayant produit
les artefacts cités, les fichiers de configuration, la matrice résultats--preuves et les
instructions de préparation des données. Lorsqu'un jeu public ne peut pas être redistribué,
la documentation doit fournir sa source, sa version attendue et la procédure de vérification.

La présence d'un chemin dans la matrice ne suffit pas si le script, l'entrée ou les
dépendances correspondantes manquent. La certification d'une reproductibilité indépendante
requiert donc une exécution par un tiers et un compte rendu des écarts. Ce contrôle externe
n'a pas été réalisé dans le cadre du présent travail.

\section{Frontière d'exploitation}

Ariel Logminer peut être exécuté localement avec un bus en mémoire et un magasin
d'idempotence SQLite, ou avec Redis lorsque plusieurs processus participent au traitement.
Le passage à une exploitation opérationnelle demanderait cependant des mécanismes non
évalués ici : gestion des secrets, contrôle d'accès, chiffrement du transport, réplication,
tests de partition réseau, sauvegarde, reprise après sinistre, supervision et politique de
rétention.

Une étude de montée en charge devrait fixer au préalable ses objectifs de latence, de
disponibilité et de reprise, employer des tâches d'inférence représentatives et injecter des
pannes du transport. Le système évalué reste un prototype de laboratoire.
